\section{Related Work}
\label{sec:related}

\paragraph{Graph unlearning and approximate updates.}
Graph unlearning asks for an already-trained GNN to remove the effect of
specified nodes, edges, or features without retraining from scratch.
Approximate GU methods are commonly grouped by mechanism: partition-based
methods retrain only affected shards in the spirit of
SISA~\cite{bourtoule2021machine} (GraphEraser~\cite{chen2022graph},
GraphRevoker); influence-function methods estimate the deletion-induced
parameter update (GIF~\cite{wu2023gif}, CEU, CGU, GST, and certified
variants such as IDEA~\cite{dong2024idea}); and learning-based methods
train deletion-aware objectives or update modules
(GNNDelete~\cite{cheng2023gnndelete}, MEGU~\cite{li2024towards}, GUKD,
D2DGN, SGU). Projection- or deployment-time methods such as Projector and
UtU form a separate category outside our scope. Existing evaluations
largely use random or protocol-driven requests; we treat the deletion set
itself as an adversarial choice and test whether vulnerability is
structured by the unlearning mechanism.

\paragraph{Adversarial deletion requests and unlearning red-teaming.}
Prior attacks on machine unlearning primarily study poisoning or
camouflaged requests in non-graph settings~\cite{marchant2022hard,di2022hidden},
where the attacker crafts samples or chooses deletions to slow down,
destabilize, or survive the unlearning procedure. Related red-team
evaluations of removal and safety mechanisms in vision models ask whether
nominally removed concepts can be recovered or bypassed~\cite{huang2024exploring,tsai2024ringabell,chin2024prompting4debugging}.
These lines motivate unlearning as an exposed safety surface. Graph
unlearning adds a distinct failure mode: deleting one node also changes
the message-passing context of its neighbors. Legal deletion requests can
therefore have non-local effects, making the central question which
deletion sets amplify approximation error beyond ordinary data removal.

\paragraph{Influence and diffusion primitives for deletion selection.}
Influence functions estimate how training examples affect predictions
~\cite{koh2017understanding}; GIF~\cite{wu2023gif} adapts this idea to
fast deletion updates, and TracIn~\cite{pruthi2020estimating} provides a
checkpoint proxy that avoids explicit Hessian inversion. Influence
maximization instead selects seed nodes with high network spread
~\cite{kempe2003maximizing}, with CELF/CELF++ accelerating greedy
selection~\cite{leskovec2007cost,goyal2011celf++}. We use these
traditions only as selection primitives: gradient-aligned sensitivity
and structural spread. Access tiers, scoring rules, and the hybrid
selector are defined in \S\ref{sec:method}.
